\documentclass[tikz,border=10pt]{standalone}
\usepackage{tikz}
\usepackage{amsmath}
\usepackage{amssymb}
\usepackage{ctex}
\usepackage{pgfplots}
\pgfplotsset{compat=1.18}
\usetikzlibrary{calc, positioning, arrows.meta, decorations.pathreplacing}

% 对称矩阵 A = [[2,1],[1,2]]
% 特征值: lambda1=3 (q1=(1,1)/sqrt2), lambda2=1 (q2=(1,-1)/sqrt2)
% 二次型 x^T A x = 1:
%   在主轴坐标 (y1,y2) 下: 3y1^2 + y2^2 = 1 → 椭圆
%   主轴方向: q1=(1,1)/sqrt2 → 半轴 1/sqrt(3), q2=(1,-1)/sqrt2 → 半轴 1
% 椭圆参数: 旋转45度，半轴 a=1/sqrt(3)≈0.577, b=1

\begin{document}
\begin{tikzpicture}[>=Stealth]

%% ====== 左图：椭圆 + 主轴方向 ======
\begin{scope}[xshift=0cm]

  \begin{axis}[
    name=ax_ellipse,
    width=7.5cm, height=7.5cm,
    xmin=-1.5, xmax=1.5,
    ymin=-1.5, ymax=1.5,
    axis lines=center,
    xlabel={$x_1$},
    ylabel={$x_2$},
    xlabel style={font=\small, right},
    ylabel style={font=\small, above},
    xtick={-1,1},
    ytick={-1,1},
    tick label style={font=\scriptsize},
    grid=both,
    grid style={gray!20},
    clip=false,
    axis equal image,
  ]

    % 椭圆 x^T A x=1，A=[[2,1],[1,2]]
    % 参数方程 (旋转45度椭圆):
    % x1 = (1/sqrt(3))*cos(t)/sqrt(2) - sin(t)/sqrt(2)  (主轴q1方向短,q2方向长)
    % 在主轴坐标: y1=q1^T x, y2=q2^T x
    % 3y1^2 + y2^2 = 1 → y1 = (1/sqrt3)cos(t), y2 = sin(t)
    % x = y1*q1 + y2*q2
    % x1 = y1/sqrt2 + y2/sqrt2 = (1/sqrt3)*cos(t)/sqrt2 + sin(t)/sqrt2
    % x2 = y1/sqrt2 - y2/sqrt2 = (1/sqrt3)*cos(t)/sqrt2 - sin(t)/sqrt2
    \addplot[blue!70!black, very thick, domain=0:360, samples=200]
      ({0.40825*cos(x) + 0.70711*sin(x)},
       {0.40825*cos(x) - 0.70711*sin(x)});
    % 标注
    \node[font=\small, blue!70!black, above right] at (axis cs:0.5,0.7)
      {$\mathbf{x}^TA\mathbf{x}=1$};

    % 主轴方向 q1=(1,1)/sqrt2 (对应lambda=3, 短轴 1/sqrt3≈0.577)
    \addplot[->,green!60!black,very thick]
      coordinates {(0,0) (0.40825,0.40825)};
    \node[font=\small, green!60!black, right] at (axis cs:0.42,0.42)
      {$\mathbf{q}_1$};
    % 主轴方向 q2=(1,-1)/sqrt2 (对应lambda=1, 长轴 1)
    \addplot[->,orange!80!black,very thick]
      coordinates {(0,0) (0.70711,-0.70711)};
    \node[font=\small, orange!80!black, below right] at (axis cs:0.70,-0.72)
      {$\mathbf{q}_2$};

    % 延长主轴虚线
    \addplot[green!40, dashed, domain=-0.6:0.6]
      ({x}, {x});
    \addplot[orange!40, dashed, domain=-1.1:1.1]
      ({x}, {-x});

    % 特征值标注
    \node[font=\scriptsize, green!50!black,
          fill=green!8, draw=green!40, rounded corners=2pt, inner sep=3pt]
      at (axis cs:-0.9,0.8) {$\lambda_1=3$（短半轴 $\tfrac{1}{\sqrt{3}}$）};
    \node[font=\scriptsize, orange!70!black,
          fill=orange!8, draw=orange!40, rounded corners=2pt, inner sep=3pt]
      at (axis cs:0.5,-1.2) {$\lambda_2=1$（长半轴 $1$）};

    % 原点
    \addplot[only marks, mark=*, mark size=1.5pt, black]
      coordinates {(0,0)};

  \end{axis}

  \node[font=\small\bfseries, align=center] at ([yshift=-0.6cm]ax_ellipse.south)
    {二次型曲线 $\mathbf{x}^TA\mathbf{x}=1$（椭圆）\\
     主轴 $=$ 特征向量方向};

\end{scope}

%% ====== 右图：谱分解公式展示 ======
\begin{scope}[xshift=9.0cm, yshift=0cm]

  % 大公式框
  \node[draw=blue!40, fill=blue!5, rounded corners=6pt,
        inner sep=16pt, align=center, font=\small]
    (decomp) at (0, 1.5) {
      \textbf{谱定理（对称矩阵）}\\[8pt]
      $A = Q\Lambda Q^T = \lambda_1\mathbf{q}_1\mathbf{q}_1^T + \lambda_2\mathbf{q}_2\mathbf{q}_2^T$\\[10pt]
      $A = \begin{pmatrix}2&1\\1&2\end{pmatrix}
         = Q\begin{pmatrix}3&0\\0&1\end{pmatrix}Q^T$\\[10pt]
      $Q = \dfrac{1}{\sqrt{2}}\begin{pmatrix}1&1\\1&-1\end{pmatrix}$\\[6pt]
      $Q^T = Q^{-1}$（正交矩阵）
    };

  % 秩一分解展开
  \node[draw=green!40, fill=green!5, rounded corners=5pt,
        inner sep=12pt, align=center, font=\small]
    at (0, -2.0) {
      \textbf{秩一展开}\\[6pt]
      $A = 3\cdot\mathbf{q}_1\mathbf{q}_1^T + 1\cdot\mathbf{q}_2\mathbf{q}_2^T$\\[4pt]
      $\mathbf{q}_1 = \dfrac{1}{\sqrt{2}}\begin{pmatrix}1\\1\end{pmatrix}$，
      $\mathbf{q}_2 = \dfrac{1}{\sqrt{2}}\begin{pmatrix}1\\-1\end{pmatrix}$\\[4pt]
      $\mathbf{q}_1\perp\mathbf{q}_2$，$\|\mathbf{q}_1\|=\|\mathbf{q}_2\|=1$
    };

  % 关键性质
  \node[draw=gray!40, fill=gray!5, rounded corners=5pt,
        inner sep=10pt, align=left, font=\small]
    at (0, -4.2) {
      \textbf{关键性质：}\\[4pt]
      $\bullet$ 实对称矩阵特征值均为\textbf{实数}\\
      $\bullet$ 不同特征值对应特征向量\textbf{正交}\\
      $\bullet$ 必可正交对角化：$A=Q\Lambda Q^T$
    };

\end{scope}

%% 整体标题
\node[font=\large\bfseries] at (4.5, 5.2)
  {谱定理：对称矩阵的正交对角化};

\end{tikzpicture}
\end{document}
